% !PRINTED 248
\weberpage{248}
Giebt es also überhaupt eine zum Exponenten $e$ gehörige Zahl $\alpha$, so giebt es so viele wie relative Primzahlen zu $e$, die positiv und kleiner als $e$ sind, eine Zahl, die wir schon früher mit $\varphi(e)$ bezeichnet haben, und die, wenn $e$ alle Divisoren von $p^n - 1$ durchläuft, der Relation

\weq{6}{\textstyle\sum \varphi(e) = p^n - 1}

genügt (Bd.\,I, \S.\,132, 133).

Ist $\psi(e)$ die Anzahl der Zahlen $\alpha$, die zum Exponenten $e$ gehören, so ist $\psi(e) = \varphi(e)$ oder $= 0$, und da die Anzahl aller Zahlen $\alpha$ gleich $p^n - 1$ ist und jede Zahl $\alpha$ zu einem und nur zu einem Exponenten gehört, so ist auch

\weq{7}{\textstyle\sum \psi(e) = p^n - 1,}

und aus (6) und (7) folgt, dass $\psi(e)$ \gesperrt{immer gleich} $\varphi(e)$ ist.

Es giebt also $\varphi(p^n - 1)$ Zahlen $\gamma$ in $\mathfrak{C}$, die zum Exponenten $p^n - 1$ gehören, und die folglich die Eigenschaft haben, dass die Potenzen

\weq{8}{1, \gamma, \gamma^2, \ldots, \gamma^{p^n-2}}

alle von einander verschieden sind. Durch diese Reihe ist daher die Gesammtheit der von Null verschiedenen Zahlen in $\mathfrak{C}$ erschöpft.

Solche Zahlen $\gamma$ heissen \gesperrt{primitive Wurzeln von $\mathfrak{C}$}.

Unter den Zahlen in $\mathfrak{C}$ giebt es solche, die als Quadrat einer anderen Zahl in $\mathfrak{C}$ darstellbar sind, die wir \gesperrt{Quadrate} nennen, und andere, bei denen dies nicht zutrifft, die wir \gesperrt{Nichtquadrate} nennen.

Wenn $p = 2$ ist, so ist jede Zahl in $\mathfrak{C}$ ein Quadrat, wie die Formel (4) zeigt.

Wenn aber $p$ ungerade ist, so besteht die eine Hälfte der von Null verschiedenen Zahlen in $\mathfrak{C}$ aus Quadraten, die andere aus Nichtquadraten.

Denn zwei entgegengesetzte Zahlen, wie $+\alpha$ und $-\alpha$, sind bei ungeradem $p$ von einander verschieden, und geben trotzdem dasselbe Quadrat. Wenn man also die sämmtlichen Zahlen von $\mathfrak{C}$ zum Quadrat erhebt, so erhält man höchstens $\frac{1}{2}(p^n - 1)$ verschiedene Quadrate. Nun kann andererseits $\beta^2$ nur dann gleich $\alpha^2$ sein, wenn $\beta = \pm\alpha$ ist; denn aus $\beta^2 - \alpha^2 = (\beta - \alpha)(\beta + \alpha) = 0$ folgt, dass $\beta - \alpha$ oder $\beta + \alpha$ verschwinden muss. Es giebt also wirklich $\frac{1}{2}(p^n - 1)$ Quadrate und ebenso viele Nichtquadrate.
